\providecommand{\LegalMathRoot}{}
\providecommand{\LegalMathBibliography}{references}
\documentclass[12pt,a4paper,oneside,openany]{book}
\usepackage[margin=25mm,headheight=15pt]{geometry}
\usepackage{fontspec}
\setmainfont{TeX Gyre Pagella}
\setsansfont{TeX Gyre Heros}
\setmonofont[Scale=0.84]{DejaVu Sans Mono}
\usepackage{amsmath,amssymb,mathtools,amsthm}
\usepackage{microtype,setspace}
\setstretch{1.06}
\usepackage{booktabs,longtable,tabularx,array}
\renewcommand{\arraystretch}{1.2}
\usepackage{graphicx,xcolor,tikz,float}
\usetikzlibrary{arrows.meta,positioning,calc,shapes.geometric}
\usepackage{enumitem,listings,titlesec}
\usepackage[round,authoryear]{natbib}
\usepackage[hyphens]{url}
\Urlmuskip=0mu plus 2mu
\definecolor{navy}{HTML}{173A53}
\definecolor{teal}{HTML}{186C75}
\definecolor{pale}{HTML}{F1F5F6}
\usepackage{xr-hyper}
\usepackage[unicode,colorlinks=true,linkcolor=navy,citecolor=navy,urlcolor=navy]{hyperref}
\hypersetup{pdftitle={LegalMath: From SFC Circulars to Reviewable Bank Controls},pdfauthor={Chak Wong},pdfsubject={SFC circular interpretation, legal argumentation, ensemble assurance and verified Java delivery}}
\externaldocument[][nocite]{\LegalMathRoot technical-companion}[technical-companion.pdf]
\titleformat{\chapter}[display]{\sffamily\bfseries\color{navy}\raggedright\hyphenpenalty=10000}{\large Chapter \thechapter}{8pt}{\huge}
\titlespacing*{\chapter}{0pt}{0pt}{24pt}
\titleformat{\section}{\Large\sffamily\bfseries\color{navy}\raggedright}{\thesection}{0.7em}{}
\titleformat{\subsection}{\large\sffamily\bfseries}{\thesubsection}{0.65em}{}
\setlength{\parindent}{0pt}
\setlength{\parskip}{0.45em}
\setlength{\emergencystretch}{3em}
\setlist{itemsep=0.15em,topsep=0.35em,leftmargin=1.6em}
\setcounter{tocdepth}{1}
\makeatletter
\def\ps@legalmath{%
  \def\@oddhead{\small\sffamily\color{navy}LegalMath\enspace|\enspace Regulatory specification and execution\hfil\leftmark}%
  \let\@evenhead\@oddhead
  \def\@oddfoot{\hfil\small\thepage\hfil}%
  \let\@evenfoot\@oddfoot}
\makeatother
\renewcommand{\chaptermark}[1]{\markboth{Chapter \thechapter}{}}
\pagestyle{legalmath}
\lstset{basicstyle=\small\ttfamily,backgroundcolor=\color{pale},frame=single,rulecolor=\color{pale},breaklines=true,columns=fullflexible,keepspaces=true,showstringspaces=false,aboveskip=0.8em,belowskip=0.8em}
\newcommand{\True}{\mathsf{T}}
\newcommand{\False}{\mathsf{F}}
\newcommand{\Unknown}{\mathsf{U}}
\newcommand{\Conflict}{\mathsf{C}}
\newcommand{\OO}{\mathsf{OUT\_OF\_SCOPE}}
\newcommand{\HKD}{\mathrm{HKD}}
\providecommand{\llbracket}{\mathopen{[\![}}
\providecommand{\rrbracket}{\mathclose{]\!]}}
\newcommand{\caseheading}[1]{\subsubsection*{#1}}
\newcommand{\IR}{\mathrm{IR}}
\newcommand{\CNL}{\mathrm{CNL}}
\newcommand{\eval}{\operatorname{eval}}
\newcommand{\srcpath}[1]{\path{#1}}
\newcolumntype{Y}{>{\raggedright\arraybackslash}X}
\newcolumntype{P}[1]{>{\raggedright\arraybackslash}p{#1}}
\newtheorem{proposition}{Proposition}[chapter]
\newtheorem{definition}{Definition}[chapter]

\input{\LegalMathRoot teaching/styles}

\begin{document}
\raggedbottom
\frontmatter
\hypersetup{pageanchor=false}
\begin{titlepage}
\vspace*{15mm}
{\sffamily\color{teal}\large REGULATORY INTERPRETATION AND BANK CONTROL\par}
\vspace{12mm}
{\sffamily\bfseries\color{navy}\fontsize{32}{37}\selectfont From SFC Circulars\\to Reviewable Bank Controls\par}
\vspace{12mm}
{\Large Interpretation, evidence and reliable execution\\for Hong Kong private banking\par}
\vspace{12mm}
{\large Chak Wong\par}
\vfill
\begin{minipage}{0.94\textwidth}
A management and research account of how a bank can turn a circular into a
reviewable control. Two circulars from Hong Kong's Securities and Futures
Commission (SFC) provide the cases. The book follows the
meaning of each requirement through interpretation, evidence, implementation
and supervised operation, and marks the point where the present evidence stops.
The technical record remains available for the people who must build and audit
the control.
\end{minipage}
\vspace{16mm}
\hrule
\vspace{5mm}
28 September 2026\hfill Draft for reader review\par
\small Public-source and synthetic examples. Independent compliance, technical
and reader acceptance remain pending. Demonstrated behaviour is distinguished
from proposed work throughout.
\end{titlepage}
\hypersetup{pageanchor=true}
\chapter*{Executive summary}
\addcontentsline{toc}{chapter}{Executive summary}
Senior management is being asked to consider a narrow question: can a bank turn
a regulatory circular into a repeatable client or product decision without
losing the meaning of the source along the way? The answer matters because a
program can apply a mistaken interpretation perfectly. The damaging error is
often made before implementation, when an exception is missed, the wrong person
is assessed, or a requirement is applied to a different activity.

LegalMath now investigates competing readings of a circular and carries the
resulting proposals into executable controls. It keeps the source passage,
interpretation, supporting facts and review decision together. A later reviewer
can reconstruct which requirement was applied, to whom, at what time and on what
evidence.

\section*{What has been done}

The work follows two public SFC circulars through the same chain of questions.
The first concerns the streamlined treatment of sophisticated professional
investors (SPIs). A synthetic client is followed through financial qualification,
product category, exposure, consent, continuing duties and a proposed
transaction. The second concerns incentives used when marketing authorised
funds. It shows why a fee discount may fall within an exception while a voucher
in the same offer may require a separate assessment. These examples are
deliberately ordinary. Their value is that a small change in a fact, a date, or
the subject being assessed changes the responsible answer.

The worked rules can be executed in Java and checked against boundary cases,
withdrawals and amendments. Later development added the investigation that comes
before those rules: generating alternative readings, looking for omitted
requirements and comparing the cases on which proposed programs disagree.
Official examples can inform that investigation, with their dates and limits
kept visible. A separate settlement-readiness circular exposed another problem:
several partial answers do not, by themselves, cover all the duties in a source.

\section*{The contribution}

LegalMath connects a circular to a bank control through an explicit, reviewable
interpretation. Compliance and legal staff state the meaning, unresolved
questions and bank-policy choices. Engineering implements that same
specification. Operations receives the decision, its evidence and the action
needed when the record is incomplete.

\newpage
The work distinguishes source duties from added bank restrictions and unknown
facts from false ones. Each fact belongs to a person, product, component, date
and activity. Preserving what was known at the time lets a correction explain
a past decision without rewriting it. A compact flag or a large language model
prompt can conceal these distinctions.

Separate duties can now be combined while rival readings remain alternatives.
A gift restriction and a disclosure duty may both need implementation; two
readings of the same restriction require a choice or an open question. The
method also looks beyond agreement. In a controlled challenge, three programs
shared a deliberately omitted exception. A separate examination of the source
exposed the gap. That is why agreement alone cannot settle the review.

\section*{What the evidence supports today}

The implementation can preserve alternatives and unfinished checks, build draft
Java controls and revisit affected work after a source change. A local amendment
exercise changed the selected control's answer while preserving the earlier
version and an unrelated control. Earlier development included live model calls;
the subsequent integration checks used retained responses and controlled cases.
Together they establish these mechanisms on public sources and synthetic facts.

One retained interpretation can now feed two executable languages: RuleIR, the
project's small rule language, and Catala, a language and compiler designed for
legal rules. Their checked common cases agree. The Catala route also executes
documented calculations over collections and partially known records that the
RuleIR target cannot express. This extends the implementation; the available
source-conversion study does not establish that either language yields better
legal interpretations or faster review.

Several decisions remain outside that evidence. The bank has not yet supplied
an approved legal interpretation for a live business perimeter, independently
validated the meanings and data mappings, measured performance against a
reference set of real cases, or connected the result to production identity,
confidentiality, audit and operational controls. The product therefore cannot
approve a client, clear a marketing offer, or replace a qualified reviewer.

\section*{The management decision}

The next decision is whether to authorize a narrow comparative study and a
supervised pilot, with a named legal owner and agreed products, entities, dates
and evidence sources. Production use would require independent legal validation,
checked data and identity mappings, security and audit controls, and evaluation
on representative cases. During shadow operation, the system's results must be
compared with qualified staff's decisions; failed or unresolved cases must remain
visible and reach a human.

\input{\LegalMathRoot frontmatter/process-map}
\tableofcontents
\mainmatter
\input{\LegalMathRoot chapters/01-problem}
\input{\LegalMathRoot chapters/02-circular}
\input{\LegalMathRoot chapters/03-proof}
\input{\LegalMathRoot chapters/04-languages}
\input{\LegalMathRoot chapters/05-search}
\input{\LegalMathRoot chapters/06-ensemble}
\input{\LegalMathRoot chapters/06a-diverse-assurance}
\input{\LegalMathRoot chapters/06b-integrated-investigation}
\input{\LegalMathRoot chapters/06c-typed-decisions}
\input{\LegalMathRoot chapters/06d-invariant-decisions}
\input{\LegalMathRoot chapters/06e-executed-assurance}
\input{\LegalMathRoot chapters/06f-input-meaning}
\input{\LegalMathRoot chapters/06g-proof-carrying-integration}
\input{\LegalMathRoot chapters/06h-successor-assurance}
\InputIfFileExists{\LegalMathRoot chapters/06h-successor-results.tex}{}{}
\input{\LegalMathRoot chapters/06i-processing-continuation}
\input{\LegalMathRoot chapters/06j-lossless-context}
\input{\LegalMathRoot chapters/06k-executable-dates}
\input{\LegalMathRoot chapters/07-verification}
\input{\LegalMathRoot chapters/08-implementation}
\input{\LegalMathRoot chapters/09-evaluation}
\input{\LegalMathRoot chapters/10-operation}
\backmatter
\begingroup
\raggedright
\bibliographystyle{plainnat}
\bibliography{\LegalMathBibliography}
\endgroup
\end{document}
